\section{Conclusions}
\label{sec:conclusions}

In this work we have computed the scattering of a guided mode at the free edge of
an n-type ZnO plate ($6mm$ symmetry) described by strain-gradient elasticity,
piezoelectricity, flexoelectricity and linearized electron drift--diffusion, with
one added carrier--strain-gradient enthalpy and its reciprocal partners.  The main
findings are as follows.

\begin{enumerate}
\item At $k_xd=0.736$ ($f=\SI{16.68}{\giga\hertz}$) the backward branch carries
\SI{1.93}{\percent} of the incident power.  Along the same branch the share rises
to \SI{94.29}{\percent} at $k_xd=0.52719$ ($f=\SI{16.6547}{\giga\hertz}$), where the
group velocity vanishes and the backward wavenumber approaches the incident one to
within \SI{0.4}{\percent}.  The response has a single peak, and its maximum is a local maximum within the window examined, and the value changes by less than $10^{-5}$
between Galerkin degrees 14 and 34.
\item These numbers are modal shares and not reflectances.  The propagating
outgoing modes carry \SI{99.98}{\percent} of the incident power at the operating
point and \SI{100.05}{\percent} at the maximum, that is, all of it to within
\SI{0.1}{\percent}; the excess over unity is of the order of the interference term, $+5.2\times10^{-4}$ at the reference plane.  At the maximum the incident flux itself has fallen by a factor
$131$ because the branch lies at a zero-group-velocity fold, so a maximal share occurs together with a minimal transported power.
\item The conversion is a property of the backward branch and not of the added
coupling.  With $H=0$ the shares at both points are those of the plate without the
coupling.  The classical plate ($\ell=0$) has a backward interval of its own, with
its fold at $k_xd=0.824$ ($f=\SI{13.69}{\giga\hertz}$); the gradient length
$\ell/d=0.25$ moves the fold to $k_xd=0.527$ and $\SI{16.65}{\giga\hertz}$.  At the
control frequencies of \cref{tab:controls}, which lie on a different classical branch, the classical plate has no backward member, the specular
reflection being its only negative-power outgoing mode.
\item The carrier coupling acts through conduction loss.  At
$|H|=0.05H^{*}$ it reduces the partition by \SI{1.21}{\percent} at the operating
point and by \SI{36.1}{\percent} at the maximum, of which 22.3 percentage points are
propagation loss between the edge and the reference plane.  When the maximum is
searched again at each coupling level it shifts and its upper bound falls, so the
coupling lowers the achievable conversion as well as displacing it.
\item The energy balance closes: the flux at the reference plane equals the edge
flux plus the conduction loss to better than $10^{-11}$ of the incident power, with
the interference term reported separately and no evanescent mode counted as
transported power.
\item At the fold the share is sensitive to loss.  It falls from $0.943$ to $0.545$ when the carrier density rises from $2\times10^{18}$ to $2\times10^{20}\,\si{\per\cubic\metre}$, and to $0.600$ and $0.197$ when a loss tangent of $10^{-5}$ and $10^{-4}$ is added to the elastic constants (\cref{sec:results_sens}); the model contains no other elastic loss, so $0.943$ is an upper bound for a plate with intrinsic damping.
\end{enumerate}

Several limitations should be kept in view.  The gradient length, the
flexoelectric coefficient and the coupling tensor are assumed values, the
carriers are linearized about a uniform doping, and the plate has a single edge.
The corner line force is small but finite in the computed field: it is
$8.4\times10^{-2}$ of the incident corner scale at the operating point and
$3.3\times10^{-3}$ at the maximum.  It decays as a power of the Galerkin degree, with the same exponent in three independent configurations, and the partition moves by
$3.9\times10^{-6}$ over the same range; the closure is therefore shown as a bound, with the decay exponent taken from the computed values.  The added
term has the structure of the chemo-mechanical coupling given in \cite{qu2025thermodynamic}, which that article neglects in its main text; the
comparison with its free-energy decomposition was made through \cite{hu2025computational} and checked against the article itself.  What the present analysis supplies is the modal partition of a slow guided branch
near its zero-group-velocity point, resolved mode by mode rather than in total power.
